\documentclass[10pt]{article}
\usepackage[T1]{fontenc}
\usepackage{lmodern,amsmath,amssymb,amsthm,mathtools,microtype}
\usepackage[margin=0.88in]{geometry}
\usepackage[colorlinks=true,urlcolor=blue,citecolor=blue,linkcolor=blue]{hyperref}
\newtheorem{theorem}{Theorem}
\newtheorem{corollary}{Corollary}
\newtheorem{remark}{Remark}
\newcommand{\cat}{\mathbin{\Vert}}
\newcommand{\ord}{\operatorname{ord}}
\title{Universal Divisibility of Fixed-Length Consecutive-Integer Concatenations\\and OEIS A075018}
\author{Qichao Wang\\{\small Hebei University of Technology, Tianjin, China}\\[0.4em]Daoyu Dong\\{\small University of Electronic Science and Technology of China, Chengdu, China}}
\date{October 17, 2026}
\begin{document}
\maketitle
\begin{abstract}
We prove that for every positive integer $N$ there is a positive integer
$k\ge3$ such that the decimal concatenation
$k\cat(k-1)\cat(k-2)\cat(k-3)$ is divisible by $N$.
Consequently, all terms of OEIS A075018 are positive.
More generally, in any even integer base, for any prescribed number $t$
of consecutive positive integers, any modulus $N$, and any specified
permutation of the $t$ blocks, one can find infinitely many
same-digit-width concatenations divisible by $N$.
The elementary constructive proof uses a base-supported/coprime
factorization of $N$, a solvable linear congruence, and the Chinese
remainder theorem. For the infinite family $N=d10^r$ with $d\mid30$
and $r\ge2$, we additionally determine the exact least value
$\mathrm{A075018}(N)=N+3$ and all admissible endpoints.
\end{abstract}

\section{The OEIS problem}
The sequence \href{https://oeis.org/A075018}{OEIS A075018}
asks for the least integer $k\ge3$ such that
\begin{equation}\label{eq:target}
 N\mid k\cat_{10}(k-1)\cat_{10}(k-2)\cat_{10}(k-3),
\end{equation}
where $N$ is the sequence index and the right-hand side is a
\emph{decimal concatenation}, not a product. The OEIS convention
assigns zero if no such $k$ exists \cite{oeis}. We show that this
zero case is impossible. The OEIS provides numerical values; a
universal nonvanishing proof is not recorded in its entry at the
time of preparation.

For positive integers $x_1,\ldots,x_t$ with exactly $d$ base-$b$
digits each, put $B=b^d$ and define
\begin{equation}\label{eq:polynomial}
 x_1\cat_b\cdots\cat_b x_t
 =\sum_{i=1}^{t}x_i B^{t-i}.
\end{equation}
Working within a single digit-length window is essential: it makes
$B$ constant, so the concatenation becomes an integer polynomial.

\section{A general even-base theorem}
\begin{theorem}[Arbitrary modulus, length, and block order]\label{thm:general}
Let $b\ge2$ be an even integer and let $N,t\ge1$ be independent
positive integers. For every bijection
$\pi:\{1,\ldots,t\}\longrightarrow\{0,\ldots,t-1\}$,
there exist infinitely many pairs of positive integers $(d,a)$ such
that the consecutive integers $a,a+1,\ldots,a+t-1$ all have
exactly $d$ base-$b$ digits and
\begin{equation}\label{eq:main}
  N\mid (a+\pi(1))\cat_b(a+\pi(2))\cat_b\cdots\cat_b(a+\pi(t)).
\end{equation}
In particular, the valid concatenated integers have unbounded length.
\end{theorem}
\begin{proof}
Factor $N=qm$, where $q$ is the largest factor of $N$ whose prime
divisors divide $b$. Then
\begin{equation}\label{eq:factor}
\gcd(q,m)=1,\qquad\gcd(b,m)=1,\qquad m\text{ is odd}.
\end{equation}
Here $q=1$ or $m=1$ is allowed, with congruences modulo $1$
understood as vacuous. Put $g=\gcd(t,m)$; because $m$ is odd,
$g$ is odd and $g\mid t$. Hence
\begin{equation}\label{eq:obstruction}
 g\mid\frac{t(t-1)}2.
\end{equation}
Indeed, $t(t-1)/2$ is an integer and twice this integer is
divisible by $g$, while $2$ is invertible modulo $g$.
It follows from the standard linear congruence criterion that
\begin{equation}\label{eq:linear}
 t a+\frac{t(t-1)}2\equiv0\pmod m
\end{equation}
has at least one solution $a_0$, and all its solutions constitute
one residue class modulo $s=m/g$.

The last block in \eqref{eq:main} is $a+\pi(t)$; impose
\begin{equation}\label{eq:q}
 a+\pi(t)\equiv0\pmod q.
\end{equation}
Since $\gcd(q,s)=1$, the Chinese remainder theorem provides
an integer $a_*$ satisfying \eqref{eq:linear} and \eqref{eq:q};
all integers $a\equiv a_*\pmod{qs}$ satisfy both.

Choose an arbitrarily large digit width $d$ with
\begin{equation}\label{eq:choose_d}
 q\mid b^d,\qquad b^d\equiv1\pmod m,\qquad
 (b-1)b^{d-1}\ge qs+t-1.
\end{equation}
These conditions can be met simultaneously: the first is automatic
for all sufficiently large $d$; for $m>1$ the second holds whenever
$d$ is a multiple of $\ord_m(b)$, and for $m=1$ it is vacuous;
the third holds eventually because $b\ge2$.
The interval of allowed first blocks
\[
 I_d=[b^{d-1},\,b^d-t]\cap\mathbb Z
\]
then contains $(b-1)b^{d-1}-t+1\ge qs$ consecutive integers,
so it contains an $a\equiv a_*\pmod{qs}$.
Thus $a,a+1,\ldots,a+t-1$ are all positive and exactly $d$ digits long.

Set $B=b^d$ and let $C$ denote the right-hand concatenation of
\eqref{eq:main}. Because $q\mid B$, the base-$B$ expansion
\eqref{eq:polynomial} gives
\[
 C\equiv a+\pi(t)\equiv0\pmod q.
\]
Because $B\equiv1\pmod m$, the order of the blocks disappears:
\[
 C\equiv\sum_{i=1}^{t}(a+\pi(i))
  =ta+\frac{t(t-1)}2\equiv0\pmod m.
\]
As $\gcd(q,m)=1$ and $N=qm$, we conclude $N\mid C$.
The admissible $d$ are arbitrarily large, producing distinct
concatenations of unbounded length $td$.
\end{proof}

\begin{remark}
The parity assumption on the base is used only to ensure that
$m$ is odd and hence \eqref{eq:obstruction} holds for every $t$.
We do not claim the analogous unrestricted result for odd bases.
The modulus $N$ and the number of blocks $t$ are \emph{independent};
this differs from the special case $N=t$ of the even-base
arithmetic-progression divisibility theorem \cite{wang}.
\end{remark}

\section{Explicit construction for A075018}
\begin{corollary}\label{cor:oeis}
For every $N\ge1$ there are infinitely many integers $k\ge3$ for
which \eqref{eq:target} holds. In particular,
\[
 \boxed{\mathrm{A075018}(N)>0\quad\text{for all }N\ge1.}
\]
\end{corollary}
\begin{proof}
Apply Theorem~\ref{thm:general} with $b=10$, $t=4$,
and $(\pi(1),\pi(2),\pi(3),\pi(4))=(3,2,1,0)$.
Then $k=a+3$, so the concatenation is precisely
$(a+3)\cat(a+2)\cat(a+1)\cat a$.
\end{proof}

Here is a particularly simple algorithm not requiring a general
linear-congruence solver. Write
\[
 N=qm,\qquad q=2^{v_2(N)}5^{v_5(N)},\qquad\gcd(m,10)=1.
\]
As $m$ is odd, $4$ is invertible modulo $m$.
The desired starting integer $a=k-3$ may be chosen by
\begin{equation}\label{eq:crt4}
 a\equiv0\pmod q,\qquad
 a\equiv-6\cdot4^{-1}\pmod m,
\end{equation}
using the Chinese remainder theorem. Let $a_*$ be its least
nonnegative representative modulo $N$. Take any sufficiently large
$d$ satisfying
\begin{equation}\label{eq:width4}
 q\mid10^d,\quad10^d\equiv1\pmod m,\quad
 9\cdot10^{d-1}\ge N+3.
\end{equation}
The least $a\ge10^{d-1}$ congruent to $a_*$ modulo $N$ is
\begin{equation}\label{eq:a_explicit}
 a=a_*+N\left\lceil\frac{10^{d-1}-a_*}{N}\right\rceil.
\end{equation}
Then $10^{d-1}\le a\le10^d-4$ and
\begin{align*}
 C&=(a+3)\cat_{10}(a+2)\cat_{10}(a+1)\cat_{10}a,\\
 C&\equiv a\equiv0\pmod q,\\
 C&\equiv4a+6\equiv0\pmod m,
\end{align*}
so $N\mid C$ and $k=a+3$ is a witness.
For example, $N=7$ admits $a=2$, giving
$C=5432=7\cdot776$ and $\mathrm{A075018}(7)=5$.
The construction need not minimize $k$.


\section{Exact least values for an infinite family}\label{sec:minima}
The general construction is not usually minimal. Nevertheless, for
a natural infinite family the minimization problem is completely
solvable, including a description of \emph{all} admissible starting
integers.

\begin{theorem}[Exact minima and complete witnesses]\label{thm:minima}
Let $r\ge2$ and let $d$ be any positive divisor of $30$,
so $d\in\{1,2,3,5,6,10,15,30\}$. Set $N=d10^r$ and
$C(k)=k\cat_{10}(k-1)\cat_{10}(k-2)\cat_{10}(k-3)$.
Then for every integer $k\ge4$,
\[
 N\mid C(k)\quad\Longleftrightarrow\quad k\equiv3\pmod N.
\]
Moreover, $C(3)=3210$ is not divisible by $N$. Thus the complete
set of admissible endpoints is
$\{N+3,2N+3,3N+3,\ldots\}$, and
\[
 \boxed{\mathrm{A075018}(d10^r)=d10^r+3
 \quad(d\mid30,\ r\ge2).}
\]
\end{theorem}
\begin{proof}
Suppose $k\ge4$ and put $x=k-3\ge1$. If $10^r\mid C(k)$,
then the last $r$ decimal digits of $C(k)$ must all vanish.
Since its final concatenated block is the ordinary decimal
representation of the positive integer $x$, this is possible
only if $x$ has at least $r$ trailing zero digits. In particular,
$x=t10^r$ for an integer $t\ge1$.

Let $\ell$ denote the number of decimal digits of $t$, and let
$U(k)=k\cat_{10}(k-1)\cat_{10}(k-2)$. Because the final block
$x=t10^r$ has $\ell+r$ digits, we have the exact identity
\[
 \frac{C(k)}{10^r}=10^\ell U(k)+t.
\]
For the prime divisors $2$ and $5$ of $d$, the first summand
is congruent to $0$. Modulo $3$, decimal concatenation is
congruent to the sum of its constituent integers:
\[
 U(k)\equiv k+(k-1)+(k-2)=3k-3\equiv0\pmod3.
\]
Since every $d\mid30$ is squarefree and has only prime
divisors in $\{2,3,5\}$, it follows that
\[
 d\mid (10^\ell U(k)+t)\quad\Longleftrightarrow\quad d\mid t.
\]
Consequently, $N=d10^r$ divides $C(k)$ exactly when
$x=k-3$ is a positive multiple of $N$. Finally,
$C(3)=3210\equiv10\pmod{100}$, whereas every $N$ in the
statement is divisible by $100$. So $k=3$ is not admissible,
and the first admissible endpoint is $N+3$.
\end{proof}

For example, $\mathrm{A075018}(100)=103$,
$\mathrm{A075018}(300)=303$,
$\mathrm{A075018}(1500)=1503$, and
$\mathrm{A075018}(3000)=3003$.
The theorem does \emph{not} assert a closed formula for
$\mathrm{A075018}(N)$ at arbitrary indices $N$.

\section{Verification and context}
The accompanying Python program checks the 65 displayed terms
of OEIS A075018 by direct search, validates constructive witnesses
for many moduli, block counts, and even bases, and checks arbitrary
permutations. These checks are independent of the universal proof.

The older OEIS A075000--A075002 work \cite{wang} deals with
related arithmetic-progression concatenations, including divisibility
by the \emph{number of blocks}. The present theorem allows an
\emph{arbitrary modulus independent of the block count} and a
prescribed ordering of the blocks. We make no claim that the
stronger general theorem is novel before exhaustive prior-art review.
Outside the exact family of Theorem~\ref{thm:minima}, determining
the least $k$ remains a separate optimization problem.

\begin{thebibliography}{9}
\bibitem{oeis}
OEIS Foundation Inc., \emph{A075018: Smallest $k$ such that the
concatenation $k,k-1,k-2,k-3$ is divisible by $n$},
\url{https://oeis.org/A075018} (accessed October 17, 2026).
\bibitem{wang}
Q. Wang, \emph{Divisibility of Concatenated Arithmetic Progressions
and Champernowne Prefixes} (2026), source and reference list at
\url{https://github.com/dongxuelian2/oeis-papers/tree/main/papers/A075000-A075002}.
\bibitem{luca}
F. Luca and B. Teguia Tabuguia, \emph{Concatenations of terms of an
arithmetic progression}, Journal of Difference Equations and Applications
\textbf{32} (2026), no.~4, 437--456,
\url{https://doi.org/10.1080/10236198.2025.2603385}.
\end{thebibliography}
\end{document}
